% !TEX root = ./../main.tex

\section{Introduction}
\label{sec:introduction}

The \acf{COP} paradigm enables writing programs that dynamically adapt their behavior in response to 
their execution context~\cite{hirschfeld08cop,appeltauer09cop,salvaneschi12survey,mens16,alegre16,cardozo22,elyasaf23}. 
\ac{COP} programs are usually organized in \emph{layers}, modular units that encapsulate behavior 
and state across different concerns, similar to the concept of \emph{object} in \ac{OOP}. Each 
\emph{layer} (also known as \emph{context} in the literature) can be activated or deactivated depending 
on the information gathered from the execution environment. Layer activation makes available to the executing program, 
the behavior and state defined within that layer. Such mechanism allows programmers to express variability, without 
having to introduce statements that toggle behavior inside the domain logic of their 
program~\cite{costanza05contextl}, achieving a better modularity and separation of concerns.

The behavior of a \ac{COP} program is determined by a set of partial methods associated with the 
layers that are active (\fref{sec:preliminaries}). A \emph{partial configuration} is 
defined as the set of layers that are active or inactive at a particular point in time of the 
execution~\cite{martou24}. As the program grows in complexity, the set of all possible partial 
configurations often suffers from a combinatorial explosion, making it harder to test and reason about 
it~\cite{martou24}, leading to bugs or execution anomalies.
Validating the correctness of a \ac{COP} program has been a topic of interest in recent years, 
and several solutions have been proposed: debuggers~\cite{uchio11}, formal 
execution models~\cite{cardozo15jist, taing16cop, cardozo18sefsas}, visualizers~\cite{duhoux18}, static 
analysis tools~\cite{cardenas23}, dynamic analysis tools~\cite{cardozo14cop}, and testing 
techniques~\cite{martou24}, among others. However, currently, no application of software verification 
techniques to \ac{COP} exists. The hypothesis of this work is that the use of software verification 
would help developers in building more reliable \ac{COP} programs, by providing 
tools to formally reason about the interaction between layers, and therefore, yield programs that are 
consistent and complete in the execution of specified adaptation behavior.

Type systems are one of the most widely used tools for guaranteeing the correctness of program 
properties. Refinement Types~\cite{jhala21} take it a step further, by allowing programmers to constraint 
a type using logical predicates, enabling more expressive and precise contracts that get automatically 
verified at compile-time. The idea is to express layer activations and interactions using 
Refinement Types (\fref{sec:state-of-the-art}), so that programmers can statically establish guarantees 
about the intended behavior of a program, without needing to dynamically validate the vast space of 
possible partial configurations. We will illustrate how Refinement Types can be used to formally verify 
\ac{COP} programs (\fref{sec:context-scala}), and discuss the advantages and drawbacks of doing 
so. The long term goal is to serve as a stepping stone to propose methodologies, design 
patterns, or even libraries that integrate software verification with \ac{COP}.

To validate our work, we built the \emph{Bellevue} graphical editor (\fref{sec:case-study}) (inspired by 
\emph{Gueze}~\cite{vallejos10pred}) as a case study. Bellevue is a sufficiently complex application to 
assess the applicability of Refinement Types in \ac{COP}, as it encompasses several operations and 
interactions between different behavioral components. We first developed the editor without Refinement Types, 
to assess and categorize the types of bugs that one would encounter in a common \ac{COP} application. 
We then adapted the editor to use Refinement Types and established some properties to detect the encountered bugs 
at compile-time.

Bellevue is written in a \ac{COP} extension for Scala using Predicated Generic 
Functions~\cite{vallejos10pred} and the Stainless framework (a verifier for higher-order functional 
programs based on Refinement Types~\cite{hamza19}). While there exist different \ac{COP} languages, 
we decided to extend Scala with \ac{COP} capabilities primarily because:
\begin{enumerate*}[label=(\arabic*)]
  \item there is already a robust verifier implemented in this language that we can reuse out of the box,
  and
  \item Scala is a higher-order functional programming language whose abstractions simplified the 
  process of creating a light-weight \ac{COP} extension
\end{enumerate*}

This work is organized as follows. \fref{sec:preliminaries} presents the preliminary concepts regarding 
\ac{COP} and type systems. \fref{sec:state-of-the-art} contains the state of the art on verification of 
\ac{COP} programs, and the main concepts to understand Refinement Types. \fref{sec:context-scala} 
defines the core abstractions of our \ac{COP} library in Scala, which is then used in 
\fref{sec:case-study} to build Bellevue and verify its layer interactions. We additionally discuss the main 
types of bugs that were found during the implementation due to inconsistent layer interactions, and how 
they can be detected at compile-time using Refinement Types. \fref{sec:related-work} presents the 
related work, showing the main verification approaches that have been found in domains similar to 
\ac{COP} (such as Self-Adaptive Systems), as well as the usage of Refinement Types to verify the 
correctness of programs in areas were \ac{COP} may be applicable (such as Cyber-Physical Systems).
Finally, \fref{sec:conclusion} concludes with the insights gained from verifying a complex \ac{COP} 
application with Refinement Types, and presents avenues of future work to improve the integration 
between \ac{COP} and Refinement Types in our library.


\endinput

That is, that all functionality associated with a context is available to complete execution of the adaptation. 
